\chapter{Mehrschichtige Perzeptrons und Backpropagation}

\section{Ziel dieses Kapitels}

In den vorherigen Kapiteln haben wir lineare Modelle betrachtet:
lineare Regression (linear regression), logistische Regression (logistic regression) und Perzeptrons (perceptrons).
Diese Modelle berechnen im Kern einen linearen Score der Form

\[
    a
    =
    \vec{w}^{\top}\vec{x}.
\]

Solche Modelle sind wichtig, aber sie haben eine zentrale Einschränkung:
Ohne geeignete Merkmalsabbildung können sie nur lineare Entscheidungsgrenzen oder lineare Zusammenhänge darstellen.

Mehrschichtige Perzeptrons (multilayer perceptrons), kurz MLPs, erweitern diese Idee.
Sie bestehen aus mehreren Schichten (layers) von künstlichen Neuronen (artificial neurons).
Durch nichtlineare Aktivierungsfunktionen (activation functions) können sie komplexe nichtlineare Funktionen lernen.

\begin{conceptbox}
Ein mehrschichtiges Perzeptron (multilayer perceptron) ist ein neuronales Netz (neural network), das Eingaben durch mehrere Schichten transformiert.

Jede Schicht berechnet eine affine Transformation und danach eine Nichtlinearität:

\[
    \vec{z}^{(l)}
    =
    g^{(l)}
    \left(
        W^{(l)}\vec{z}^{(l-1)}+\vec{b}^{(l)}
    \right).
\]
\end{conceptbox}

Nach diesem Kapitel solltest du folgende Fragen beantworten können:

\begin{itemize}
    \item Was ist ein mehrschichtiges Perzeptron (multilayer perceptron)?
    \item Was ist der Unterschied zwischen Eingabeschicht, versteckter Schicht und Ausgabeschicht?
    \item Warum braucht man nichtlineare Aktivierungsfunktionen?
    \item Was ist ein Forward Pass?
    \item Was ist eine Verlustfunktion (loss function)?
    \item Was ist Backpropagation?
    \item Wie berechnet man Gradienten mit der Kettenregel?
    \item Was bedeuten lokale Gradienten und Fehlerterme?
    \item Wie sehen die Gradienten für Gewichte und Biases aus?
    \item Wie trainiert man ein neuronales Netz mit Gradientenabstieg?
    \item Warum können Gradienten verschwinden oder explodieren?
    \item Welche Rolle spielen Sigmoid, Tanh und ReLU?
\end{itemize}

\section{Vom Perzeptron zum neuronalen Netz}

Ein einzelnes Perzeptron berechnet

\[
    a
    =
    \vec{w}^{\top}\vec{x}+b
\]

und danach eine Aktivierungsfunktion:

\[
    z
    =
    g(a).
\]

Beim klassischen Perzeptron ist \(g\) eine Schwellenfunktion.
In modernen neuronalen Netzen verwendet man meistens differenzierbare Aktivierungsfunktionen wie Sigmoid, Tanh oder ReLU.

\begin{definitionbox}
Ein künstliches Neuron (artificial neuron) berechnet

\[
    a
    =
    \vec{w}^{\top}\vec{x}+b,
\]

und danach

\[
    z
    =
    g(a).
\]

Dabei ist \(a\) die Voraktivierung (pre-activation), \(z\) die Aktivierung (activation), \(\vec{w}\) der Gewichtsvektor, \(b\) der Bias und \(g\) die Aktivierungsfunktion.
\end{definitionbox}

Ein einzelnes Neuron ist noch stark eingeschränkt.
Die Idee eines neuronalen Netzes ist, viele solcher Neuronen in Schichten zu organisieren.

\section{Schichten eines neuronalen Netzes}

Ein mehrschichtiges Perzeptron besteht aus mehreren Schichten:

\begin{itemize}
    \item Eingabeschicht (input layer),
    \item eine oder mehrere versteckte Schichten (hidden layers),
    \item Ausgabeschicht (output layer).
\end{itemize}

Die Eingabeschicht enthält die Eingabedaten.
Die versteckten Schichten berechnen interne Repräsentationen.
Die Ausgabeschicht erzeugt die Vorhersage.

\begin{definitionbox}
Eine versteckte Schicht (hidden layer) ist eine Schicht zwischen Eingabe und Ausgabe.

Sie heißt versteckt, weil ihre Zielwerte nicht direkt im Datensatz beobachtet werden.
Sie wird nur indirekt durch das Trainingsziel gelernt.
\end{definitionbox}

Für ein Netz mit \(L\) Schichten schreiben wir:

\[
    \vec{z}^{(0)}
    =
    \vec{x}.
\]

Die Eingabe \(\vec{x}\) ist also die Aktivierung der Schicht \(0\).
Für Schicht \(l\) gilt dann:

\[
    \vec{a}^{(l)}
    =
    W^{(l)}\vec{z}^{(l-1)}
    +
    \vec{b}^{(l)},
\]

und

\[
    \vec{z}^{(l)}
    =
    g^{(l)}(\vec{a}^{(l)}).
\]

Dabei ist:

\[
    l \in \{1,\dots,L\}.
\]

Die Ausgabe des Netzes ist:

\[
    \vec{y}
    =
    \vec{z}^{(L)}.
\]

\section{Dimensionen der Parameter}

Angenommen, Schicht \(l-1\) hat \(D_{l-1}\) Einheiten und Schicht \(l\) hat \(D_l\) Einheiten.
Dann gilt:

\[
    \vec{z}^{(l-1)} \in \mathbb{R}^{D_{l-1}},
\]

\[
    W^{(l)} \in \mathbb{R}^{D_l \times D_{l-1}},
\]

\[
    \vec{b}^{(l)} \in \mathbb{R}^{D_l},
\]

\[
    \vec{a}^{(l)} \in \mathbb{R}^{D_l},
\]

und

\[
    \vec{z}^{(l)} \in \mathbb{R}^{D_l}.
\]

\begin{notebox}
Die Gewichtsmatrix \(W^{(l)}\) hat eine Zeile pro Neuron der aktuellen Schicht und eine Spalte pro Ausgabe der vorherigen Schicht.

Das Produkt

\[
    W^{(l)}\vec{z}^{(l-1)}
\]

hat daher die Dimension \(\mathbb{R}^{D_l}\).
\end{notebox}

\section{Forward Pass}

Der Forward Pass ist die Berechnung der Ausgabe aus der Eingabe.

Für ein Netz mit zwei Schichten, also einer versteckten Schicht und einer Ausgabeschicht, gilt:

\[
    \vec{z}^{(0)}
    =
    \vec{x}.
\]

Die versteckte Schicht berechnet:

\[
    \vec{a}^{(1)}
    =
    W^{(1)}\vec{x}
    +
    \vec{b}^{(1)},
\]

\[
    \vec{z}^{(1)}
    =
    g^{(1)}(\vec{a}^{(1)}).
\]

Die Ausgabeschicht berechnet:

\[
    \vec{a}^{(2)}
    =
    W^{(2)}\vec{z}^{(1)}
    +
    \vec{b}^{(2)},
\]

\[
    \vec{y}
    =
    \vec{z}^{(2)}
    =
    g^{(2)}(\vec{a}^{(2)}).
\]

\begin{definitionbox}
Der Forward Pass ist die schichtweise Berechnung

\[
    \vec{x}
    \longrightarrow
    \vec{z}^{(1)}
    \longrightarrow
    \vec{z}^{(2)}
    \longrightarrow
    \dots
    \longrightarrow
    \vec{y}.
\]

Dabei werden die aktuellen Parameter \(W^{(l)}\) und \(\vec{b}^{(l)}\) verwendet.
\end{definitionbox}

\section{Warum Nichtlinearitäten notwendig sind}

Angenommen, ein Netz hätte zwei lineare Schichten ohne Aktivierungsfunktion:

\[
    \vec{z}^{(1)}
    =
    W^{(1)}\vec{x}
    +
    \vec{b}^{(1)}
\]

und

\[
    \vec{y}
    =
    W^{(2)}\vec{z}^{(1)}
    +
    \vec{b}^{(2)}.
\]

Einsetzen ergibt:

\[
    \vec{y}
    =
    W^{(2)}
    \left(
        W^{(1)}\vec{x}
        +
        \vec{b}^{(1)}
    \right)
    +
    \vec{b}^{(2)}.
\]

Also:

\[
    \vec{y}
    =
    W^{(2)}W^{(1)}\vec{x}
    +
    W^{(2)}\vec{b}^{(1)}
    +
    \vec{b}^{(2)}.
\]

Setze:

\[
    W
    =
    W^{(2)}W^{(1)}
\]

und

\[
    \vec{b}
    =
    W^{(2)}\vec{b}^{(1)}
    +
    \vec{b}^{(2)}.
\]

Dann:

\[
    \vec{y}
    =
    W\vec{x}
    +
    \vec{b}.
\]

Das ist wieder nur eine affine Abbildung.

\begin{conceptbox}
Mehrere lineare Schichten ohne Nichtlinearität sind äquivalent zu einer einzigen linearen Schicht.

Nichtlineare Aktivierungsfunktionen sind deshalb notwendig, damit tiefe Netze komplexe nichtlineare Funktionen darstellen können.
\end{conceptbox}

\section{Wichtige Aktivierungsfunktionen}

\subsection{Sigmoid}

Die Sigmoid-Funktion (sigmoid function) ist:

\[
    \sigma(a)
    =
    \frac{1}{1+\exp(-a)}.
\]

Sie bildet auf das Intervall \((0,1)\) ab.

Die Ableitung ist:

\[
    \sigma'(a)
    =
    \sigma(a)(1-\sigma(a)).
\]

\begin{notebox}
Sigmoid eignet sich als Ausgabeaktivierung für binäre Klassifikation.
In versteckten Schichten kann Sigmoid problematisch sein, weil die Ableitung für große positive oder negative Werte sehr klein wird.
\end{notebox}

\subsection{Tanh}

Die Tanh-Funktion ist:

\[
    \tanh(a)
    =
    \frac{\exp(a)-\exp(-a)}
    {\exp(a)+\exp(-a)}.
\]

Sie bildet auf das Intervall \((-1,1)\) ab.

Die Ableitung ist:

\[
    \frac{d}{da}\tanh(a)
    =
    1-\tanh^2(a).
\]

\subsection{ReLU}

Die ReLU-Funktion (rectified linear unit) ist:

\[
    \operatorname{ReLU}(a)
    =
    \max(0,a).
\]

Also:

\[
    \operatorname{ReLU}(a)
    =
    \begin{cases}
        a, & a>0, \\
        0, & a\leq 0.
    \end{cases}
\]

Die Ableitung ist für \(a\neq 0\):

\[
    \operatorname{ReLU}'(a)
    =
    \begin{cases}
        1, & a>0, \\
        0, & a<0.
    \end{cases}
\]

\begin{conceptbox}
ReLU ist in modernen neuronalen Netzen sehr beliebt, weil sie einfach ist und für positive Aktivierungen keinen kleinen Sättigungsgradienten hat.
\end{conceptbox}

\section{Ausgabeschichten}

Die Wahl der Ausgabeschicht hängt vom Problemtyp ab.

\subsection{Regression}

Für Regression verwendet man oft eine lineare Ausgabe:

\[
    \vec{y}
    =
    \vec{a}^{(L)}.
\]

Bei eindimensionaler Regression:

\[
    y \in \mathbb{R}.
\]

Eine typische Verlustfunktion ist der quadratische Fehler:

\[
    E
    =
    \frac{1}{2}
    (y-t)^2.
\]

\subsection{Binäre Klassifikation}

Für binäre Klassifikation verwendet man meistens eine Sigmoid-Ausgabe:

\[
    y
    =
    \sigma(a^{(L)}).
\]

Dann gilt:

\[
    y
    =
    p(t=1\mid\vec{x}).
\]

Die passende Verlustfunktion ist die binäre Kreuzentropie (binary cross-entropy):

\[
    E
    =
    -\left[
        t\log y
        +
        (1-t)\log(1-y)
    \right].
\]

\subsection{Mehrklassen-Klassifikation}

Für \(K\) Klassen verwendet man Softmax:

\[
    y_k
    =
    \frac{\exp(a_k)}
    {\sum_{j=1}^{K}\exp(a_j)}.
\]

Die passende Verlustfunktion ist die Mehrklassen-Kreuzentropie (multiclass cross-entropy):

\[
    E
    =
    -\sum_{k=1}^{K}
    t_k\log y_k.
\]

\section{Verlustfunktion für einen Datensatz}

Für einen einzelnen Datenpunkt schreiben wir den Verlust als

\[
    E_n(\Theta),
\]

wobei \(\Theta\) alle Parameter des Netzes bezeichnet:

\[
    \Theta
    =
    \{W^{(1)},\vec{b}^{(1)},\dots,W^{(L)},\vec{b}^{(L)}\}.
\]

Für einen Datensatz mit \(N\) Datenpunkten ist der Gesamtverlust häufig:

\[
    E(\Theta)
    =
    \sum_{n=1}^{N}
    E_n(\Theta).
\]

Oder als Durchschnitt:

\[
    E(\Theta)
    =
    \frac{1}{N}
    \sum_{n=1}^{N}
    E_n(\Theta).
\]

\begin{definitionbox}
Training eines neuronalen Netzes bedeutet, die Parameter \(\Theta\) so zu wählen, dass eine Verlustfunktion minimiert wird:

\[
    \Theta^*
    =
    \arg\min_{\Theta}
    E(\Theta).
\]
\end{definitionbox}

\section{Das Problem: Viele Parameter}

Ein neuronales Netz kann sehr viele Parameter haben.
Für jede Schicht \(l\) gibt es:

\[
    W^{(l)}
\]

und

\[
    \vec{b}^{(l)}.
\]

Um Gradientenabstieg zu verwenden, brauchen wir alle Ableitungen:

\[
    \frac{\partial E}{\partial W^{(l)}}
\]

und

\[
    \frac{\partial E}{\partial \vec{b}^{(l)}}.
\]

Direkt alle Ableitungen einzeln auszurechnen wäre unübersichtlich und ineffizient.
Backpropagation löst genau dieses Problem.

\begin{conceptbox}
Backpropagation ist ein effizienter Algorithmus zur Berechnung aller Gradienten in einem neuronalen Netz.

Die zentrale mathematische Grundlage ist die Kettenregel (chain rule).
\end{conceptbox}

\section{Kettenregel}

Backpropagation basiert auf der Kettenregel.
Für skalare Funktionen gilt:

\[
    z = f(y),
    \qquad
    y = g(x).
\]

Dann ist:

\[
    \frac{dz}{dx}
    =
    \frac{dz}{dy}
    \frac{dy}{dx}.
\]

In neuronalen Netzen sind die Größen meistens Vektoren und Matrizen.
Die Idee bleibt aber gleich:
Eine Änderung in einer frühen Schicht beeinflusst alle späteren Schichten.

\begin{definitionbox}
Die Kettenregel erlaubt, Ableitungen zusammengesetzter Funktionen zu berechnen.

Für neuronale Netze bedeutet das:
Der Fehler wird von der Ausgabe zurück durch die Schichten propagiert.
\end{definitionbox}

\section{Backpropagation-Idee}

Der Forward Pass berechnet:

\[
    \vec{x}
    \longrightarrow
    \vec{z}^{(1)}
    \longrightarrow
    \dots
    \longrightarrow
    \vec{z}^{(L)}
    =
    \vec{y}
    \longrightarrow
    E.
\]

Backpropagation läuft in umgekehrter Richtung:

\[
    E
    \longrightarrow
    \vec{z}^{(L)}
    \longrightarrow
    \vec{a}^{(L)}
    \longrightarrow
    \dots
    \longrightarrow
    \vec{a}^{(1)}.
\]

Dabei werden sogenannte Fehlerterme (error terms) berechnet.

\begin{definitionbox}
Der Fehlerterm einer Schicht \(l\) ist

\[
    \vec{\delta}^{(l)}
    =
    \frac{\partial E}{\partial \vec{a}^{(l)}}.
\]

Er beschreibt, wie stark der Verlust auf Änderungen der Voraktivierung \(\vec{a}^{(l)}\) reagiert.
\end{definitionbox}

\section{Ein einfaches Zwei-Schichten-Netz}

Wir betrachten nun ein Netz mit einer versteckten Schicht und einer Ausgabeschicht.

Forward Pass:

\[
    \vec{z}^{(0)}
    =
    \vec{x}.
\]

Versteckte Schicht:

\[
    \vec{a}^{(1)}
    =
    W^{(1)}\vec{x}
    +
    \vec{b}^{(1)},
\]

\[
    \vec{z}^{(1)}
    =
    g^{(1)}(\vec{a}^{(1)}).
\]

Ausgabeschicht:

\[
    \vec{a}^{(2)}
    =
    W^{(2)}\vec{z}^{(1)}
    +
    \vec{b}^{(2)},
\]

\[
    \vec{y}
    =
    \vec{z}^{(2)}
    =
    g^{(2)}(\vec{a}^{(2)}).
\]

Der Verlust ist:

\[
    E
    =
    E(\vec{y},\vec{t}).
\]

\section{Gradienten der Ausgabeschicht}

Wir wollen die Ableitung nach \(W^{(2)}\) berechnen.
Für die Ausgabeschicht gilt:

\[
    \vec{a}^{(2)}
    =
    W^{(2)}\vec{z}^{(1)}
    +
    \vec{b}^{(2)}.
\]

Der Fehlerterm ist:

\[
    \vec{\delta}^{(2)}
    =
    \frac{\partial E}{\partial \vec{a}^{(2)}}.
\]

Dann gilt für die Gewichtsmatrix:

\[
    \frac{\partial E}{\partial W^{(2)}}
    =
    \vec{\delta}^{(2)}
    (\vec{z}^{(1)})^{\top}.
\]

Für den Bias gilt:

\[
    \frac{\partial E}{\partial \vec{b}^{(2)}}
    =
    \vec{\delta}^{(2)}.
\]

\begin{definitionbox}
Für eine Schicht \(l\) gilt allgemein:

\[
    \frac{\partial E}{\partial W^{(l)}}
    =
    \vec{\delta}^{(l)}
    (\vec{z}^{(l-1)})^{\top},
\]

und

\[
    \frac{\partial E}{\partial \vec{b}^{(l)}}
    =
    \vec{\delta}^{(l)}.
\]
\end{definitionbox}

\section{Fehlerterm der Ausgabeschicht}

Der Fehlerterm der Ausgabeschicht hängt von der gewählten Ausgabefunktion und Verlustfunktion ab.

\subsection{Quadratischer Fehler}

Bei quadratischem Fehler:

\[
    E
    =
    \frac{1}{2}
    \|\vec{y}-\vec{t}\|^2
\]

gilt:

\[
    \frac{\partial E}{\partial \vec{y}}
    =
    \vec{y}-\vec{t}.
\]

Wenn

\[
    \vec{y}
    =
    g^{(L)}(\vec{a}^{(L)}),
\]

dann ist:

\[
    \vec{\delta}^{(L)}
    =
    (\vec{y}-\vec{t})
    \odot
    (g^{(L)})'(\vec{a}^{(L)}).
\]

Dabei bedeutet \(\odot\) komponentenweise Multiplikation.

\subsection{Softmax mit Kreuzentropie}

Bei Softmax-Ausgabe und Kreuzentropie vereinfacht sich der Fehlerterm zu:

\[
    \vec{\delta}^{(L)}
    =
    \vec{y}-\vec{t}.
\]

\subsection{Sigmoid mit binärer Kreuzentropie}

Bei Sigmoid-Ausgabe und binärer Kreuzentropie gilt ebenfalls:

\[
    \delta^{(L)}
    =
    y-t.
\]

\begin{conceptbox}
Die Kombination aus Sigmoid und Kreuzentropie oder Softmax und Kreuzentropie führt zu einem einfachen Ausgabefehler:

\[
    \text{Ausgabefehler}
    =
    \text{Vorhersage}
    -
    \text{Ziel}.
\]
\end{conceptbox}

\section{Fehlerterm der versteckten Schicht}

Nun wollen wir den Fehlerterm der versteckten Schicht berechnen:

\[
    \vec{\delta}^{(1)}
    =
    \frac{\partial E}{\partial \vec{a}^{(1)}}.
\]

Die versteckte Schicht beeinflusst den Verlust über die Ausgabeschicht.
Daher verwenden wir die Kettenregel.

Für allgemeine Schicht \(l\) gilt:

\[
    \vec{\delta}^{(l)}
    =
    \left(
        (W^{(l+1)})^{\top}
        \vec{\delta}^{(l+1)}
    \right)
    \odot
    (g^{(l)})'(\vec{a}^{(l)}).
\]

Für die versteckte Schicht im Zwei-Schichten-Netz:

\[
    \vec{\delta}^{(1)}
    =
    \left(
        (W^{(2)})^{\top}
        \vec{\delta}^{(2)}
    \right)
    \odot
    (g^{(1)})'(\vec{a}^{(1)}).
\]

\begin{definitionbox}
Die Backpropagation-Rekursion für versteckte Schichten lautet:

\[
    \vec{\delta}^{(l)}
    =
    \left(
        (W^{(l+1)})^{\top}
        \vec{\delta}^{(l+1)}
    \right)
    \odot
    (g^{(l)})'(\vec{a}^{(l)}).
\]
\end{definitionbox}

\section{Backpropagation-Algorithmus}

Für einen einzelnen Datenpunkt läuft Backpropagation so ab:

\begin{enumerate}
    \item Forward Pass:
    Berechne für alle Schichten
    \[
        \vec{a}^{(l)}
        =
        W^{(l)}\vec{z}^{(l-1)}
        +
        \vec{b}^{(l)}
    \]
    und
    \[
        \vec{z}^{(l)}
        =
        g^{(l)}(\vec{a}^{(l)}).
    \]

    \item Verlust berechnen:
    \[
        E
        =
        E(\vec{z}^{(L)},\vec{t}).
    \]

    \item Ausgabefehler berechnen:
    \[
        \vec{\delta}^{(L)}
        =
        \frac{\partial E}{\partial \vec{a}^{(L)}}.
    \]

    \item Rückwärts durch die versteckten Schichten:
    \[
        \vec{\delta}^{(l)}
        =
        \left(
            (W^{(l+1)})^{\top}
            \vec{\delta}^{(l+1)}
        \right)
        \odot
        (g^{(l)})'(\vec{a}^{(l)}).
    \]

    \item Gradienten berechnen:
    \[
        \frac{\partial E}{\partial W^{(l)}}
        =
        \vec{\delta}^{(l)}
        (\vec{z}^{(l-1)})^{\top},
    \]
    \[
        \frac{\partial E}{\partial \vec{b}^{(l)}}
        =
        \vec{\delta}^{(l)}.
    \]
\end{enumerate}

\begin{conceptbox}
Backpropagation speichert im Forward Pass die Aktivierungen und Voraktivierungen.
Im Backward Pass werden daraus effizient alle Gradienten berechnet.
\end{conceptbox}

\section{Gradientenabstieg für MLPs}

Nach der Berechnung der Gradienten werden die Parameter aktualisiert.
Für Batch-Gradientenabstieg gilt:

\[
    W^{(l)}
    \leftarrow
    W^{(l)}
    -
    \eta
    \frac{\partial E}{\partial W^{(l)}},
\]

und

\[
    \vec{b}^{(l)}
    \leftarrow
    \vec{b}^{(l)}
    -
    \eta
    \frac{\partial E}{\partial \vec{b}^{(l)}}.
\]

Dabei ist \(\eta>0\) die Lernrate (learning rate).

\begin{definitionbox}
Training mit Gradientenabstieg bedeutet:

\[
    \Theta
    \leftarrow
    \Theta
    -
    \eta
    \nabla_{\Theta}E(\Theta).
\]

Dabei enthält \(\Theta\) alle Gewichte und Biases des Netzes.
\end{definitionbox}

\section{Mini-Batch-Training}

In der Praxis verwendet man selten den ganzen Datensatz auf einmal.
Stattdessen verwendet man Mini-Batches.

Ein Mini-Batch \(\mathcal{B}\) ist eine Teilmenge der Trainingsdaten.
Der Mini-Batch-Verlust ist:

\[
    E_{\mathcal{B}}(\Theta)
    =
    \frac{1}{|\mathcal{B}|}
    \sum_{n\in\mathcal{B}}
    E_n(\Theta).
\]

Die Parameter werden mit dem Gradienten dieses Mini-Batch-Verlusts aktualisiert:

\[
    \Theta
    \leftarrow
    \Theta
    -
    \eta
    \nabla_{\Theta}E_{\mathcal{B}}(\Theta).
\]

\begin{conceptbox}
Mini-Batch-Training ist ein Kompromiss zwischen Batch-Gradientenabstieg und stochastischem Gradientenabstieg.

Es ist effizient, skaliert gut und nutzt moderne Hardware besser.
\end{conceptbox}

\section{Matrixform für mehrere Datenpunkte}

Für einen Mini-Batch kann man die Aktivierungen als Matrizen schreiben.
Sei \(Z^{(0)}\) eine Matrix, deren Spalten die Eingaben des Mini-Batches sind.

Für Schicht \(l\):

\[
    A^{(l)}
    =
    W^{(l)}Z^{(l-1)}
    +
    \vec{b}^{(l)}\vec{1}^{\top}.
\]

Die Aktivierungen sind:

\[
    Z^{(l)}
    =
    g^{(l)}(A^{(l)}).
\]

Die Fehlerterme werden ebenfalls als Matrizen geschrieben:

\[
    \Delta^{(l)}
    =
    \frac{\partial E}{\partial A^{(l)}}.
\]

Dann gilt:

\[
    \frac{\partial E}{\partial W^{(l)}}
    =
    \Delta^{(l)}
    (Z^{(l-1)})^{\top},
\]

und

\[
    \frac{\partial E}{\partial \vec{b}^{(l)}}
    =
    \Delta^{(l)}\vec{1}.
\]

\begin{notebox}
Je nach Konvention liegen Datenpunkte als Zeilen oder Spalten in den Matrizen.
Die Formeln müssen dann entsprechend transponiert werden.
Wichtig ist nicht die Konvention, sondern konsistente Dimensionen.
\end{notebox}

\section{Initialisierung der Gewichte}

Die Wahl der Anfangswerte ist wichtig.
Wenn alle Gewichte mit \(0\) initialisiert werden, sind Neuronen derselben Schicht symmetrisch.
Dann lernen sie dieselben Features.

\begin{conceptbox}
Gewichte neuronaler Netze sollten zufällig initialisiert werden, damit verschiedene Neuronen unterschiedliche Funktionen lernen können.
\end{conceptbox}

Typische Initialisierungen berücksichtigen die Anzahl der Eingänge und Ausgänge einer Schicht.

\subsection{Xavier-Initialisierung}

Xavier-Initialisierung wird häufig für Sigmoid oder Tanh verwendet.
Die Varianz der Gewichte wird so gewählt, dass Aktivierungen und Gradienten nicht zu stark wachsen oder schrumpfen.

\subsection{He-Initialisierung}

He-Initialisierung wird häufig für ReLU-Netze verwendet.
Sie berücksichtigt, dass ReLU ungefähr die Hälfte der Aktivierungen auf \(0\) setzt.

\begin{notebox}
Die genaue Initialisierung ist ein praktisches Detail, aber die Grundidee ist wichtig:
Zu große Gewichte können Aktivierungen und Gradienten explodieren lassen.
Zu kleine Gewichte können Signale verschwinden lassen.
\end{notebox}

\section{Verschwindende und explodierende Gradienten}

Bei tiefen Netzen werden viele Ableitungen miteinander multipliziert.
Dadurch können Gradienten sehr klein oder sehr groß werden.

Wenn viele Faktoren kleiner als \(1\) sind, kann der Gradient verschwinden:

\[
    \|\nabla E\|
    \approx
    0.
\]

Das nennt man verschwindende Gradienten (vanishing gradients).

Wenn viele Faktoren größer als \(1\) sind, kann der Gradient explodieren:

\[
    \|\nabla E\|
    \gg
    1.
\]

Das nennt man explodierende Gradienten (exploding gradients).

\begin{conceptbox}
Verschwindende Gradienten erschweren das Lernen früher Schichten.

Explodierende Gradienten machen das Training instabil.
\end{conceptbox}

Sigmoid und Tanh können verschwindende Gradienten begünstigen, wenn sie in Sättigungsbereiche geraten.
ReLU reduziert dieses Problem teilweise, hat aber eigene Probleme wie tote Neuronen.

\section{Tote ReLU-Neuronen}

Ein ReLU-Neuron ist tot, wenn es für alle oder fast alle Eingaben

\[
    a \leq 0
\]

hat.
Dann ist:

\[
    \operatorname{ReLU}(a)=0
\]

und

\[
    \operatorname{ReLU}'(a)=0.
\]

Der Gradient für dieses Neuron ist dann \(0\), und es kann kaum noch lernen.

\begin{notebox}
Tote ReLUs können durch zu große Lernraten oder ungünstige Initialisierung entstehen.
Varianten wie Leaky ReLU reduzieren dieses Problem.
\end{notebox}

\section{Regularisierung neuronaler Netze}

Neuronale Netze können sehr flexibel sein und daher überanpassen.
Typische Regularisierungsmethoden sind:

\begin{itemize}
    \item \(L_2\)-Regularisierung,
    \item Early Stopping,
    \item Dropout,
    \item Datenaugmentation (data augmentation),
    \item kleinere Netzwerkarchitekturen.
\end{itemize}

\subsection{\(L_2\)-Regularisierung}

Man ergänzt den Verlust um:

\[
    \frac{\lambda}{2}
    \sum_l
    \|W^{(l)}\|_F^2.
\]

Dann ist:

\[
    E_{\mathrm{reg}}(\Theta)
    =
    E(\Theta)
    +
    \frac{\lambda}{2}
    \sum_l
    \|W^{(l)}\|_F^2.
\]

\subsection{Early Stopping}

Beim Early Stopping wird das Training beendet, wenn der Validierungsfehler nicht mehr besser wird.

\begin{conceptbox}
Early Stopping wirkt als Regularisierung, weil das Netz nicht beliebig lange an die Trainingsdaten angepasst wird.
\end{conceptbox}

\subsection{Dropout}

Bei Dropout werden während des Trainings zufällig Aktivierungen auf \(0\) gesetzt.
Dadurch soll das Netz robuster werden und sich weniger auf einzelne Neuronen verlassen.

\section{Universal Approximation}

Ein wichtiger theoretischer Satz sagt:
Ein neuronales Netz mit mindestens einer versteckten Schicht und geeigneter Aktivierungsfunktion kann unter bestimmten Bedingungen jede stetige Funktion auf einem kompakten Bereich beliebig gut approximieren.

\begin{conceptbox}
Der Universal-Approximation-Satz sagt nicht, dass ein Netz automatisch gut lernt.

Er sagt nur, dass es prinzipiell genug Ausdruckskraft haben kann.
Training, Datenmenge, Optimierung und Regularisierung bleiben trotzdem entscheidend.
\end{conceptbox}

\section{Vergleich mit linearen Modellen}

\[
\begin{array}{c|c|c}
\text{Modell}
&
\text{Form}
&
\text{Nichtlinearität}
\\
\hline
\text{Lineare Regression}
&
\vec{w}^{\top}\vec{x}
&
\text{keine}
\\
\hline
\text{Logistische Regression}
&
\sigma(\vec{w}^{\top}\vec{x})
&
\text{nur Ausgabe}
\\
\hline
\text{MLP}
&
g^{(L)}(W^{(L)}\cdots g^{(1)}(W^{(1)}\vec{x}+\vec{b}^{(1)}))
&
\text{mehrere Schichten}
\end{array}
\]

\begin{conceptbox}
Lineare und logistische Regression lernen eine direkte Abbildung aus den gegebenen Features.

Ein MLP lernt zusätzlich interne Features in den versteckten Schichten.
\end{conceptbox}

\section{Häufige Fehler}

\begin{examtipbox}
Typische Fehler bei MLPs und Backpropagation:

\begin{enumerate}
    \item Denken, mehrere lineare Schichten ohne Aktivierungsfunktion seien mächtiger als eine lineare Schicht.
    \item Die Dimensionen von \(W^{(l)}\), \(\vec{z}^{(l)}\) und \(\vec{b}^{(l)}\) verwechseln.
    \item Voraktivierung \(\vec{a}^{(l)}\) und Aktivierung \(\vec{z}^{(l)}\) verwechseln.
    \item Die Ableitung der Aktivierungsfunktion in Backpropagation vergessen.
    \item Den Fehlerterm \(\vec{\delta}^{(l)}\) mit dem Gradienten nach den Gewichten verwechseln.
    \item Bei \(\frac{\partial E}{\partial W^{(l)}}\) das äußere Produkt vergessen.
    \item Softmax mit Sigmoid verwechseln.
    \item Kreuzentropie und quadratischen Fehler unpassend kombinieren.
    \item Gewichte mit \(0\) initialisieren und dadurch Symmetrieprobleme erzeugen.
    \item Training Error und Validation Error verwechseln.
\end{enumerate}
\end{examtipbox}

\section{Mini-Aufgaben}

\subsection{Aufgabe 1: Dimensionen einer Schicht}

Eine Schicht erhält Eingaben

\[
    \vec{z}^{(l-1)}\in\mathbb{R}^{5}
\]

und hat

\[
    3
\]

Neuronen.
Welche Dimensionen haben \(W^{(l)}\), \(\vec{b}^{(l)}\), \(\vec{a}^{(l)}\) und \(\vec{z}^{(l)}\)?

\begin{examplebox}
Da die Schicht \(3\) Ausgaben und \(5\) Eingaben hat:

\[
    W^{(l)}\in\mathbb{R}^{3\times 5}.
\]

Der Bias hat eine Komponente pro Neuron:

\[
    \vec{b}^{(l)}\in\mathbb{R}^{3}.
\]

Daher:

\[
    \vec{a}^{(l)}\in\mathbb{R}^{3},
    \qquad
    \vec{z}^{(l)}\in\mathbb{R}^{3}.
\]
\end{examplebox}

\subsection{Aufgabe 2: ReLU berechnen}

Berechne

\[
    \operatorname{ReLU}(-2),
    \qquad
    \operatorname{ReLU}(3).
\]

\begin{examplebox}
Mit

\[
    \operatorname{ReLU}(a)=\max(0,a)
\]

folgt:

\[
    \operatorname{ReLU}(-2)=0,
\]

und

\[
    \operatorname{ReLU}(3)=3.
\]
\end{examplebox}

\subsection{Aufgabe 3: Ausgabefehler bei Softmax}

Für einen Datenpunkt sei

\[
    \vec{y}
    =
    \begin{pmatrix}
        0{,}1 \\
        0{,}7 \\
        0{,}2
    \end{pmatrix},
    \qquad
    \vec{t}
    =
    \begin{pmatrix}
        0 \\
        1 \\
        0
    \end{pmatrix}.
\]

Berechne den Fehlerterm der Ausgabeschicht bei Softmax mit Kreuzentropie.

\begin{examplebox}
Bei Softmax mit Kreuzentropie gilt:

\[
    \vec{\delta}^{(L)}
    =
    \vec{y}-\vec{t}.
\]

Also:

\[
    \vec{\delta}^{(L)}
    =
    \begin{pmatrix}
        0{,}1 \\
        0{,}7 \\
        0{,}2
    \end{pmatrix}
    -
    \begin{pmatrix}
        0 \\
        1 \\
        0
    \end{pmatrix}
    =
    \begin{pmatrix}
        0{,}1 \\
        -0{,}3 \\
        0{,}2
    \end{pmatrix}.
\]
\end{examplebox}

\subsection{Aufgabe 4: Gewichtgradient einer Schicht}

Gegeben seien

\[
    \vec{\delta}^{(l)}
    =
    \begin{pmatrix}
        2 \\
        -1
    \end{pmatrix},
    \qquad
    \vec{z}^{(l-1)}
    =
    \begin{pmatrix}
        3 \\
        4
    \end{pmatrix}.
\]

Berechne

\[
    \frac{\partial E}{\partial W^{(l)}}.
\]

\begin{examplebox}
Es gilt:

\[
    \frac{\partial E}{\partial W^{(l)}}
    =
    \vec{\delta}^{(l)}
    (\vec{z}^{(l-1)})^{\top}.
\]

Also:

\[
    \frac{\partial E}{\partial W^{(l)}}
    =
    \begin{pmatrix}
        2 \\
        -1
    \end{pmatrix}
    \begin{pmatrix}
        3 & 4
    \end{pmatrix}
    =
    \begin{pmatrix}
        6 & 8 \\
        -3 & -4
    \end{pmatrix}.
\]
\end{examplebox}

\subsection{Aufgabe 5: Warum Nichtlinearität?}

Warum ist ein Netz aus mehreren linearen Schichten ohne Aktivierungsfunktionen nicht mächtiger als eine einzelne lineare Schicht?

\begin{examplebox}
Zwei lineare Schichten ergeben:

\[
    \vec{h}=W^{(1)}\vec{x},
\]

\[
    \vec{y}=W^{(2)}\vec{h}.
\]

Einsetzen liefert:

\[
    \vec{y}=W^{(2)}W^{(1)}\vec{x}.
\]

Setze

\[
    W=W^{(2)}W^{(1)}.
\]

Dann:

\[
    \vec{y}=W\vec{x}.
\]

Das ist wieder eine einzige lineare Schicht.
\end{examplebox}

\section{Zusammenfassung}

\begin{itemize}
    \item Ein mehrschichtiges Perzeptron (multilayer perceptron) besteht aus mehreren Schichten künstlicher Neuronen.
    \item Eine Schicht berechnet:
    \[
        \vec{a}^{(l)}
        =
        W^{(l)}\vec{z}^{(l-1)}
        +
        \vec{b}^{(l)}.
    \]
    \item Danach folgt eine Aktivierung:
    \[
        \vec{z}^{(l)}
        =
        g^{(l)}(\vec{a}^{(l)}).
    \]
    \item Nichtlineare Aktivierungsfunktionen sind notwendig, weil mehrere lineare Schichten sonst wieder nur eine lineare Abbildung ergeben.
    \item Der Forward Pass berechnet die Ausgabe des Netzes aus der Eingabe.
    \item Backpropagation berechnet Gradienten effizient mit der Kettenregel.
    \item Der Fehlerterm einer Schicht ist:
    \[
        \vec{\delta}^{(l)}
        =
        \frac{\partial E}{\partial \vec{a}^{(l)}}.
    \]
    \item Für versteckte Schichten gilt:
    \[
        \vec{\delta}^{(l)}
        =
        \left(
            (W^{(l+1)})^{\top}
            \vec{\delta}^{(l+1)}
        \right)
        \odot
        (g^{(l)})'(\vec{a}^{(l)}).
    \]
    \item Die Gradienten sind:
    \[
        \frac{\partial E}{\partial W^{(l)}}
        =
        \vec{\delta}^{(l)}
        (\vec{z}^{(l-1)})^{\top},
    \]
    \[
        \frac{\partial E}{\partial \vec{b}^{(l)}}
        =
        \vec{\delta}^{(l)}.
    \]
    \item Bei Softmax mit Kreuzentropie gilt:
    \[
        \vec{\delta}^{(L)}
        =
        \vec{y}-\vec{t}.
    \]
    \item Bei Sigmoid mit binärer Kreuzentropie gilt:
    \[
        \delta^{(L)}
        =
        y-t.
    \]
    \item MLPs werden meist mit Mini-Batch-Gradientenabstieg trainiert.
    \item Initialisierung, Aktivierungsfunktionen und Regularisierung sind entscheidend für stabiles Training.
\end{itemize}

\section{Typische Prüfungsfragen}

\begin{examtipbox}
Mögliche Prüfungsfragen zu diesem Kapitel:

\begin{enumerate}
    \item Was ist ein mehrschichtiges Perzeptron?
    \item Was ist der Unterschied zwischen Eingabeschicht, versteckter Schicht und Ausgabeschicht?
    \item Warum braucht man nichtlineare Aktivierungsfunktionen?
    \item Zeige, dass mehrere lineare Schichten ohne Aktivierung äquivalent zu einer linearen Schicht sind.
    \item Was ist ein Forward Pass?
    \item Was ist Backpropagation?
    \item Welche Rolle spielt die Kettenregel?
    \item Definiere den Fehlerterm \(\vec{\delta}^{(l)}\).
    \item Leite die Backpropagation-Rekursion für versteckte Schichten her.
    \item Wie berechnet man \(\frac{\partial E}{\partial W^{(l)}}\)?
    \item Wie berechnet man \(\frac{\partial E}{\partial \vec{b}^{(l)}}\)?
    \item Warum vereinfacht sich der Ausgabefehler bei Softmax mit Kreuzentropie zu \(\vec{y}-\vec{t}\)?
    \item Was ist der Unterschied zwischen Sigmoid, Tanh und ReLU?
    \item Was sind verschwindende Gradienten?
    \item Was sind explodierende Gradienten?
    \item Warum sollte man Gewichte nicht alle mit \(0\) initialisieren?
    \item Was ist Mini-Batch-Training?
    \item Was ist Early Stopping?
    \item Was ist Dropout?
    \item Was sagt der Universal-Approximation-Satz aus und was sagt er nicht aus?
\end{enumerate}
\end{examtipbox}